\bookchapter{第十二章：影音娛樂與閲讀}

\chapterintro{手機而家基本上已經取代咗MP3、電視遙控器、Kindle同掌上遊戲機。但係講真，好多人淨係用咗呢啲App最基本嘅功能，仲有好多好用嘅地方未發現。呢一章我哋就嚟傾下點樣等手機娛樂更順手、更舒服。}

\booksection{一、聽音樂：揀邊個App，點樣聽先舒服}

音樂App而家功能都差唔多，但各有各嘅側重點，揀啱可以省唔少心：

\begin{itemize}
  \item \textbf{網易雲音樂}：佢嘅強項係歌單社區同評論區。你想聽「適合落雨天聽嘅歌」或者「學到攰時放鬆用嘅純音樂」，搜一下就有別人整理好嘅歌單。評論區有時候比歌仲精彩。缺點係部分歌曲版權不全，有啲歌聽到聽到就灰咗。
  \item \textbf{QQ音樂}：版權最全，周杰倫、五月天呢啲熱門歌手嘅基本上淨係QQ音樂有。音質選項都好豐富，會員可以開無損音質。如果你對「想聽嘅歌都搜得到」呢件事好在意，QQ音樂係最穩妥嘅選擇。
  \item \textbf{Apple Music}：蘋果用戶嘅首選。介面乾淨、冇廣告、音質好，而且支援無損同杜比全景聲。但佢嘅推薦算法不如國內App識中文用戶嘅喜好，歌單文化都冇網易雲咁發達。適合鍾意簡潔體驗、唔太在乎社交功能嘅人。
  \item \textbf{Spotify}：全球用戶最多，推薦算法公認最強。「每週發現」同「日推」準確率高得嚇人，用得越耐越識你。缺點係免費版廣告多，而且某些中文歌版權覆蓋不全。需要特殊網絡先用得到。
\end{itemize}

\textbf{幾個聽歌小技巧}：如果你用嘅係藍牙耳機，去音樂App設定入面睇下有冇「音效調節」或「均衡器」，稍微調一下低音同高音，效果提升明顯。另外，串流媒體聽歌好食流量，記得喺設定入面將「僅WiFi下播放無損」打開，或者提前下載好鍾意嘅歌單離線聽。

\booksection{二、播客（Podcast）：通勤路上嘅「耳朵伴侶」}

播客呢樣嘢，未聽過嘅人覺得同自己無關，一旦開始聽就返唔到轉頭。本質就係「電台節目嘅Netflix版」——有人傾科技，有人講歷史，有人讀故事，你想聽乜嘢類型嘅都有。

\begin{itemize}
  \item \textbf{小宇宙App}：國內做得最好嘅播客App，介面清爽，冇亂七八糟嘅功能。佢嘅「每日推薦」同「編輯精選」幫你發現新節目，評論區都好活躍。我個人聽得最多嘅係《忽左忽右》（文化歷史類）同《不合時宜》（社會議題類），每期都好似喺旁聽幾個聰明人傾偈。
  \item \textbf{蘋果播客}：iPhone自帶，唔需要額外下載App。同CarPlay嘅整合最好，揸車嗰陣好方便。
  \item \textbf{收聽技巧}：大部分播客App都支援1.2倍速或1.5倍速播放，習慣咗之後可以省唔少時間。我一般1.2倍聽中文，1.5倍聽英文。另外一定要開「自動下載新節目」同「播完自動刪除」，咁樣你嘅手機唔會堆滿聽過嘅音頻檔案。
\end{itemize}

\booksection{三、睇影片：由播放器到串流媒體嘅正確打開方式}

好多人唔知，手機上可以播嘅格式遠比串流媒體App入面提供嘅多。一個好嘅本地播放器可以解決好多問題。

\begin{itemize}
  \item \textbf{VLC播放器}：開源免費，幾乎乜嘢格式都播到——MKV、AVI、FLAC、甚至係ISO光碟鏡像。如果你下載咗一個影片喺電腦上播到但手機上播唔到，八成係格式問題，裝個VLC就搞掂。佢仲支援外掛字幕，將.srt字幕檔案同影片放喺同一個資料夾，VLC會自動載入。
  \item \textbf{nPlayer}：如果你想要更精緻嘅體驗，呢個付費App值得買。介面靚，手勢控制好順手——左邊上下滑調亮度，右邊上下滑調音量，左右滑快進快退。支援DTS同杜比音效，睇大片嗰陣臨場感強好多。
  \item \textbf{Bilibili（B站）}已經唔淨係二次元網站，而家上面嘅學習資源多到離譜。考研、編程、攝影、烹飪，乜嘢教程都有。彈幕文化係B站嘅靈魂，睇影片嗰陣開住彈幕，感覺好似有一羣人喺陪你一齊睇。不過睇學習影片嗰陣建議關彈幕，唔係容易分心。
  \item \textbf{投屏到電視}：手機螢幕再大都唔如電視睇得舒服。同一個WiFi下，打開影片App右上角嘅「TV」投屏掣，就可以將畫面投射到電視上。大部分智能電視同投影儀都支援，睇劇睇戲體驗好太多。iPhone用戶仲可以用AirPlay直接鏡像整個螢幕。
\end{itemize}

\booksection{四、電子書閲讀：手機上都可以好好睇書}

有人話手機上讀書傷眼、容易分心，呢啲確實有道理。但手機閲讀嘅優勢都好明顯——隨時隨地，碎片時間利用最大化。我喺地鐵上、排隊嗰陣、瞓前，用手機讀咗唔少書。

\begin{itemize}
  \item \textbf{微信讀書}：目前國內體驗最好嘅閲讀App。書庫好全，出版書同網文都有。佢用「閲讀時長換書幣」嘅模式，睇得多就可以免費換書，對學生黨好友好。佢嘅「劃線」功能可以睇別人喺邊啲段落做咗標記，有時候睇到別人標註嘅好句子，確實有共鳴。不過「閲讀排行榜」呢樣嘢有啲卷，容易變成「為咗刷時長而讀書」，建議關掉。
  \item \textbf{Kindle App}：如果你有Kindle設備，手機端Kindle App可以同步閲讀進度。Kindle嘅正版電子書排版質量比盜版PDF好太多，但Kindle中國區書城已經關咗，需要切換到海外區。
  \item \textbf{Apple Books（圖書）}：iPhone自帶，適合讀EPUB同PDF。介面極簡，翻頁動畫好舒服。如果你由網上下載咗電子書檔案，直接「分享到圖書」就可以匯入。
  \item \textbf{閲讀設定建議}：唔好用邊個App，一定要將背景設成暖黃色調（護眼模式），字體調大啲，行距拉開。夜晚讀書打開「深色模式」——黑底白字，眼舒服好多。千萬唔好白底黑字細字體喺暗光下讀太耐，眼疲勞係累積嘅。
\end{itemize}

\booksection{五、手機遊戲：玩得爽，都唔傷手機}

我對手機遊戲嘅態度係：玩可以，但唔好等佢反過嚟控制你。合理使用遊戲工具，玩得更舒服，都唔耽誤正事。

\begin{itemize}
  \item \textbf{遊戲模式/遊戲空間}：而家主流手機都有呢個功能。打開之後手機會自動清後台、遮蔽通知、防止誤觸——你正在打團戰，突然一個電話彈出嚟，呢種嘢遊戲模式可以幫你擋掉。部分手機仲支援性能模式/E-sports模式，強制處理器跑滿，但代價係發熱同耗電都快啲。
  \item \textbf{畫質同幀率點樣調}：唔好無腦拉滿。先保證幀率（流暢度），再追求畫質。60幀比30幀舒服太多，但120幀對手機壓力好大，中端機唔使強求。如果你玩嘅係競技類遊戲（比如王者、吃雞），幀率優先級遠高於畫質——流暢嘅操作手感比好睇嘅草皮重要一百倍。
  \item \textbf{遊戲時手機發燙點算}：呢個係正常嘅，但長期高温會損傷電池壽命。幾個降温辦法：除低手機殼散熱；將螢幕亮度降到70\%左右（亮度係高發熱源）；唔好一邊充電一邊打機（發熱疊加，電池死得快）。如果手機燙到手握唔住嘅程度，趕緊退出嚟等佢歇一陣。
  \item \textbf{防沉迷小工具}：iOS嘅「螢幕使用時間」同安卓嘅「數碼健康」都可以畀遊戲App設每日時間限額。我畀自己設嘅係工作日最多一個鐘，週末兩個鐘。到時間咗App圖示會變灰，撳唔入去。呢個限制係真係有用，唔係裝樣嘅。
  \item \textbf{雲遊戲}：如果你手機性能唔夠但係又想玩大作，可以試下雲遊戲平臺（比如騰訊START雲遊戲、網易雲遊戲）。遊戲喺遠端伺服器上跑，畫面以影片流嘅方式傳到你嘅手機上。對網絡要求高但對你手機性能幾乎冇要求。缺點係延遲感不可避免，適合玩策略類、RPG類，唔適合FPS射擊類。
\end{itemize}

\begin{tcbpractice}{嚟動手：打造你嘅手機娛樂系統}
\begin{enumerate}
  \item 喺音樂App入面下載你最鍾意嘅3張專輯或歌單，確保離線可聽。
  \item 去小宇宙或蘋果播客訂閲3個感興趣嘅頻道，下載一期試聽。
  \item 安裝VLC播放器，嘗試用佢播放一段電腦拷貝過嚟嘅影片檔案。
  \item 喺微信讀書入面揾一本想讀嘅書，設定暖色背景同護眼模式，將字體調到適合自己嘅大細。
  \item 如果你打機，打開遊戲模式睇下有邊啲設定可以優化（免打擾、幀率優先等）。
\end{enumerate}

\end{tcbpractice}
